\section{Related Work}

\subsection{LLM Calibration}

Neural network calibration has been extensively studied since \citet{guo2017calibration} established Expected Calibration Error (ECE) as the standard metric and introduced temperature scaling as a post-hoc calibration method. For LLMs specifically, \citet{kadavath2022language} demonstrated that models possess intrinsic self-knowledge---they can predict whether their own answers are correct at above-chance rates. This suggests that calibration information exists within models but may not surface in standard prompting.

Subsequent work explored various approaches to elicit this calibration information. \citet{lin2022teaching} trained models to express uncertainty in natural language, while \citet{xiong2023can} provided a comprehensive evaluation of confidence elicitation methods, comparing direct prompting, multi-sample consistency, and linguistic confidence markers. \citet{tian2023just} specifically studied prompting strategies for calibration, demonstrating that ``just asking'' for confidence can yield reasonable estimates.

\subsection{Chain-of-Thought Prompting}

\citet{wei2022cot} introduced chain-of-thought prompting, showing that prompting models to reason step-by-step dramatically improves performance on reasoning tasks. \citet{kojima2022zeroshot} extended this to zero-shot settings with the simple prompt ``Let's think step by step.'' \citet{wang2023selfconsistency} introduced self-consistency, using multiple samples with majority voting to improve reliability.

These works focus primarily on accuracy rather than calibration. A key observation from this literature is that CoT produces visible reasoning traces, which could in principle reveal model uncertainty through linguistic markers like hedging words and qualifications.

\subsection{Uncertainty Quantification in NLP}

The NLP literature has long recognized linguistic hedging as markers of epistemic uncertainty \citep{hyland1998hedging}. In the context of LLMs, \citet{kuhn2023semantic} proposed semantic uncertainty, clustering semantically equivalent responses to estimate uncertainty. This approach leverages response diversity as an implicit uncertainty signal.

Our work differs from prior approaches in systematically isolating the mechanism by which CoT might improve calibration. While \citet{tian2023just} tested various prompting strategies and \citet{xiong2023can} evaluated elicitation methods, neither decomposed the CoT+confidence combination into its component mechanisms. We provide the first systematic verification that (1) CoT produces hedging markers, (2) these markers are structurally positioned to influence confidence generation, and (3) they negatively correlate with verbalized confidence, supporting a ``self-reading'' interpretation.
